Este capítulo cubre 




\section{Protección de datos y privacidad (RGPD)}

\section{Licencias de software y herramientas de terceros}

\section{Uso de la información del sector público (Open data)}

\section{Propiedad intelectual del proyecto}

\section{Ética profesional}



\textit{ \textbf{Posible opción}
El desarrollo de este Trabajo de Fin de Grado se ha guiado por los principios fundamentales del Código de Ética y
Práctica Profesional de Ingeniería de Software (ACM/IEEE). Desde una perspectiva deontológica, la ingeniería 
informática no consiste únicamente en la resolución técnica de problemas, sino en la aplicación del conocimiento 
empírico para orientar la tecnología hacia el bien común. En este contexto, el simulador diseñado busca un beneficio 
social tangible e inmediato: proporcionar a la Empresa Municipal de Transportes una herramienta objetiva para 
optimizar las dotaciones programadas, mejorando así la eficiencia y sostenibilidad de la movilidad urbana. Si bien 
es inherente a la ingeniería de software que los sistemas complejos pueden generar consecuencias operativas o sociales 
imprevisibles, la responsabilidad ética del diseño se ha salvaguardado mediante la aplicación de medidas rigurosas y 
estándares contrastados de la profesión. El uso de validaciones estadísticas exhaustivas y metodologías robustas como 
la simulación de Monte Carlo garantiza los más altos niveles de calidad técnica y fiabilidad, limitando los riesgos 
derivados de la puesta en producción del software.
Asimismo, el compromiso ético del proyecto trasciende el estricto cumplimiento legal impuesto por normativas de
privacidad como el RGPD. En el ámbito de la automatización y el diseño algorítmico, se ha prestado especial atención a
la mitigación de sesgos computacionales que pudieran derivar en discriminación o en un análisis deficiente de la
calidad del servicio. Por ello, el sistema modela el comportamiento de la red basándose de forma estricta en datos
operacionales empíricos y agregados. Al descartar explícitamente dotar a los agentes simulados de capacidades
deliberativas arbitrarias sin fundamento en los datos, se evita la introducción de sesgos de confirmación o decisiones
algorítmicas opacas. De este modo, se asegura que la herramienta evalúe el rendimiento del transporte público de manera
transparente y equitativa, respetando el interés general de la ciudadanía y evitando cualquier discriminación en la
planificación de los recursos urbanos.}